\chapter{Creio na remissão dos pecados}
\chaptermark{Creio na remissão dos pecados}

Este capítulo aprofunda o terceiro artigo do Símbolo Apostólico, professando a fé no perdão dos pecados. Desdobramos os parágrafos correspondentes do Catecismo da Igreja Católica (CIC 976--987), destacando o Batismo como primeiro sacramento de perdão e a Penitência como reconciliação para os batizados. Para adultos em processo de iniciação cristã pelo RICA, essa fé impulsiona a confissão sacramental, unindo catecúmenos à misericórdia pascal.

\section{O Batismo, Primeiro Sacramento do Perdão (\S 976--980)}

O Credo associa a fé no perdão dos pecados à fé no Espírito Santo, na Igreja e na comunhão dos santos; Cristo confere aos apóstolos o poder divino de perdoar pela efusão do Espírito: ``Recebei o Espírito Santo. Àqueles a quem perdoardes os pecados, ser-lhes-ão perdoados'' (Jo 20,23).

Nosso Senhor liga o perdão à fé e ao Batismo: ``Quem crer e for batizado será salvo'' (Mc 16,16); é o primeiro e principal sacramento de perdão, unindo-nos a Cristo morto e ressuscitado para nova vida (Rm 6,4).

No Batismo, o perdão é pleno e completo: nada resta a apagar, nem pecado original, nem culpas pessoais, nem penas de expiação; a graça batismal liberta totalmente, mas não da fraqueza da natureza.

Na luta contra a concupiscência, ninguém escapa a toda ferida do pecado; a Igreja, com as chaves do Reino, deve perdoar penitentes até o último instante, além do Batismo.

Pela sacramento da Penitência, o batizado reconcilia-se com Deus e a Igreja, chamado "batismo laborioso" pelos Padres, necessário para salvação como o Batismo para os não regenerados.

Para adultos a partir dos 20 anos no RICA, o Batismo no rito de iniciação pascal perdoa integralmente, mas prepara para Penitência quaresmal, purificando no catecumenato rumo à Eucaristia.

\section{O Ministério da Reconciliação na Igreja (\S 981--987)}

Após a Ressurreição, Cristo envia apóstolos para pregar penitência e perdão em seu nome a todas as nações (Lc 24,47); exercem o ``ministério da reconciliação'' por Batismo e chaves, revivendo almas pelo sangue de Cristo e Espírito.

Nenhum pecado é irreparável: a Igreja perdoa todo ofensor arrependido; Cristo deseja portas abertas a quem se converte (Mt 18,21-22).

A catequese desperta fé no dom de Cristo à Igreja: poder de perdoar via apóstolos e sucessores, confirmando sacerdotes como instrumentos divinos, sem os quais não haveria esperança eterna.

O Credo liga ``perdão dos pecados'' ao Espírito Santo, pois o Ressuscitado confiou esse poder aos apóstolos ao dar-lhes o Espírito.

Batismo é primeiro sacramento de perdão, unindo a Cristo e dando o Espírito; sacerdotes e sacramentos são instrumentos de Cristo para justificar.

Na catequese para adultos não iniciados, o ministério da reconciliação no tempo pascal da mistagogia fortalece neófitos na confissão regular, integrando misericórdia à comunhão eclesial.

\section{Conclusão}

Professar o perdão dos pecados afirma Batismo e Penitência como dons trinitários na Igreja. Adultos no RICA, abracem essa graça sacramental para vida nova em Cristo.